\begin{conclusion}[全谱的纤维分裂刚性：固定 Kronecker 体谱簇与 \texorpdfstring{$(q+1)$}{(q+1)} 维例外块]\label{con:xi-terminal-replica-softcore-fiber-splitting-full-spectrum}
固定整数 \(q\ge 2\)，沿用 \(U_q=\{0,1\}^q\)、\(D^{(q)}=K^{\otimes q}\) 与
\(
T^{(q)}=\frac12(\mathbf{1}\mathbf{1}^{\mathsf T}+D^{(q)})
\)
的记号（\S\ref{sec:pom}）。
定义固定谱簇数列
\[
\nu_k^{(q)}:=\frac12(-1)^k\varphi^{q-2k}\qquad(0\le k\le q),
\]
则 \(T^{(q)}\) 的谱由两部分完全刻画：
\begin{enumerate}
  \item \textbf{（固定谱簇）}
  对每个 \(0\le k\le q\)，数 \(\nu_k^{(q)}\) 作为 \(T^{(q)}\) 的特征值出现，其重数恰为 \(\binom{q}{k}-1\)；
  并约定 \(k=0,q\) 时该重数为 \(0\)（因 \(\binom{q}{0}=\binom{q}{q}=1\)）。
  \item \textbf{（例外谱）}
  除上述固定谱簇外，\(T^{(q)}\) 还恰有 \(q+1\) 个互异的实特征值；
  它们与 \(S_q\)-对称降维矩阵 \(A^{(q)}\) 的谱完全一致（定义 \eqref{eq:pom-multiplicity-composition-Aq-def}，命题 \ref{prop:pom-multiplicity-composition-sq-symmetry-reduction}），并且全部为单根。
\end{enumerate}
\end{conclusion}
\begin{proof}[证明要点]
由 \(D^{(q)}=K^{\otimes q}\) 的张量本征分解，\(D^{(q)}\) 的特征值为 \((-1)^k\varphi^{q-2k}\)（重数 \(\binom{q}{k}\)）。
并且 \(\mathbf{1}\mathbf{1}^{\mathsf T}\) 为秩一算子：若 \(x\perp\mathbf{1}\)，则 \((\mathbf{1}\mathbf{1}^{\mathsf T})x=0\)，从而 \(T^{(q)}x=\frac12D^{(q)}x\)。
\par
另一方面，结论 \ref{con:xi-terminal-replica-softcore-exceptional-secular-stieltjes} 的权重计算给出 \(P_k\mathbf{1}\neq 0\)（\(0\le k\le q\)），其中 \(P_k\) 为 \(D^{(q)}\) 对应第 \(k\) 个谱层的正交投影。
故每个谱层与 \(\mathbf{1}^\perp\) 的交维数为 \(\binom{q}{k}-1\)，得到固定谱簇的重数。
剩余维数为
\(
2^q-\sum_{k=0}^{q}(\binom{q}{k}-1)=q+1
\)，并由 \(S_q\)-对称子空间承载；
其限制矩阵即为 \(A^{(q)}\)（命题 \ref{prop:pom-multiplicity-composition-sq-symmetry-reduction}），从而例外谱为 \(\mathrm{spec}(A^{(q)})\)。
单根性由同一结论的“对角 \(+\) 秩一”世俗刻画与权重全非零推出。证毕。
\end{proof}

\begin{conclusion}[例外谱的世俗方程：双代数权重的 Stieltjes 型刻画]\label{con:xi-terminal-replica-softcore-exceptional-secular-stieltjes}
沿用结论 \ref{con:xi-terminal-replica-softcore-fiber-splitting-full-spectrum} 的记号。
令 \(K\) 的单位正交本征向量为 \((e_+,e_-)\)，对应本征值 \((\varphi,-\varphi^{-1})\)。
将 \(u:=(1,1)^{\mathsf T}\) 写为
\[
u=\alpha e_+ + \beta e_-,
\qquad
\alpha^2=1+\frac{2}{\sqrt5},\quad \beta^2=1-\frac{2}{\sqrt5}.
\]
对 \(0\le k\le q\) 定义权重
\[
w_k^{(q)}:=\binom{q}{k}\alpha^{2(q-k)}\beta^{2k}
=\binom{q}{k}\Bigl(1+\frac{2}{\sqrt5}\Bigr)^{q-k}\Bigl(1-\frac{2}{\sqrt5}\Bigr)^{k},
\]
并仍记固定谱簇对角值 \(\nu_k^{(q)}=\frac12(-1)^k\varphi^{q-2k}\)。
则例外谱 \(\{\lambda_i^{(q)}\}_{i=0}^{q}=\mathrm{spec}(A^{(q)})\) 恰为下列世俗方程的全部（且全为单根）实解：
\[
\boxed{\ \frac12\sum_{k=0}^{q}\frac{w_k^{(q)}}{\lambda-\nu_k^{(q)}}=1\ }.
\]
等价地，令
\[
P_q(\lambda)
:=\prod_{k=0}^{q}(\lambda-\nu_k^{(q)})
-\frac12\sum_{k=0}^{q}w_k^{(q)}\prod_{\substack{0\le j\le q\\ j\ne k}}(\lambda-\nu_j^{(q)}),
\]
则 \(\{\lambda_i^{(q)}\}\) 为 \(P_q(\lambda)=0\) 的根集；
并且 \(P_q(\lambda)\in\QQ[\lambda]\)，更强地 \(2^{q+1}P_q(\lambda)\in\ZZ[\lambda]\)。
\end{conclusion}
\begin{proof}
由 \(\mathbf{1}=u^{\otimes q}\) 与 \(D^{(q)}=K^{\otimes q}\) 的张量本征分解，
\[
\mathbf{1}
=\sum_{k=0}^{q}\alpha^{q-k}\beta^k\sum_{|S|=k}e_-^{\otimes S}\otimes e_+^{\otimes S^{\mathsf c}},
\]
其中对固定 \(k\) 的内层向量两两正交且位于 \(D^{(q)}\) 的第 \(k\) 个谱层。
故 \(P_k\mathbf{1}\neq 0\) 且
\[
\|P_k\mathbf{1}\|^2
=\binom{q}{k}\alpha^{2(q-k)}\beta^{2k}
=w_k^{(q)}.
\]
令 \(f_k:=P_k\mathbf{1}/\|P_k\mathbf{1}\|\) 并记 \(c_k:=\|P_k\mathbf{1}\|=\sqrt{w_k^{(q)}}\)。
在正交基 \(\{f_k\}_{k=0}^{q}\) 下，限制到 \(\mathrm{span}\{P_k\mathbf{1}\}_{k=0}^{q}\) 的算子为
\[
T^{(q)}\big|_{\mathrm{sym}}
=\mathrm{diag}(\nu_0^{(q)},\dots,\nu_q^{(q)})+\frac12\,cc^{\mathsf T}.
\]
对“对角 \(+\) 秩一”矩阵应用行列式引理得到
\[
0=\det\bigl(\lambda I-\mathrm{diag}(\nu_k^{(q)})-\tfrac12cc^{\mathsf T}\bigr)
=\det(\lambda I-\mathrm{diag}(\nu_k^{(q)}))
\Bigl(1-\frac12\sum_{k=0}^{q}\frac{c_k^2}{\lambda-\nu_k^{(q)}}\Bigr),
\]
从而得到所示世俗方程，并且由于 \(c_k^2=w_k^{(q)}>0\)，各根均为单根且与 \(\{\nu_k^{(q)}\}\) 严格交错。
\par
多项式表示由两边同乘 \(\prod_{k}(\lambda-\nu_k^{(q)})\) 得到。
至于有理/整性：例外谱亦等于 \(A^{(q)}\) 的谱（命题 \ref{prop:pom-multiplicity-composition-sq-symmetry-reduction}），而 \(A^{(q)}\) 在基 \(\{e_b\}_{b=0}^{q}\) 下为半整数矩阵（\eqref{eq:pom-multiplicity-composition-Aq-def}），故其特征多项式属于 \(\QQ[\lambda]\)，并且乘以 \(2^{q+1}\) 后系数皆为整数。证毕。
\end{proof}

\begin{conclusion}[对称降维算子的显式逆核与整性]\label{con:xi-terminal-replica-softcore-Aq-explicit-inverse-integrality}
固定整数 \(q\ge 2\)，并沿用 \(S_q\)-对称降维矩阵 \(A^{(q)}\) 的定义 \eqref{eq:pom-multiplicity-composition-Aq-def}。
则 \(A^{(q)}\) 可逆，且其逆矩阵具有闭式核：对任意 \(0\le a,b\le q\)，有
\[
\boxed{\ (A^{(q)})^{-1}_{a,b}
=2(-1)^{a+b-q}\binom{a}{q-b}-\delta_{a,0}\delta_{b,0}\ } ,
\]
其中约定当 \(q-b>a\) 时 \(\binom{a}{q-b}=0\)。
特别地，\(A^{(q)}\) 的逆矩阵逐项为整数。
\end{conclusion}
\begin{proof}
令
\[
L_{a,b}:=\binom{q-a}{b},\qquad
R_{a,b}:=\binom{q}{b},
\qquad
\Rightarrow\quad
A^{(q)}=\frac12(L+R).
\]
注意 \(R=uv^{\mathsf T}\) 为秩一矩阵，其中 \(u\) 为全 \(1\) 列向量、\(v_b=\binom{q}{b}\)。
对 \(L+uv^{\mathsf T}\) 应用 Sherman--Morrison 公式并用 \(v^{\mathsf T}L^{-1}u=1\)（二项反演的等价表述），得到
\[
(L+uv^{\mathsf T})^{-1}=L^{-1}-\frac12\,L^{-1}uv^{\mathsf T}L^{-1}.
\]
而 \(L^{-1}u=e_0\) 且 \(v^{\mathsf T}L^{-1}=e_0^{\mathsf T}\)，故
\[
(L+uv^{\mathsf T})^{-1}=L^{-1}-\frac12\,e_0e_0^{\mathsf T},
\qquad
(A^{(q)})^{-1}=2(L+uv^{\mathsf T})^{-1}=2L^{-1}-e_0e_0^{\mathsf T}.
\]
最后，\(L\) 经行反转与 Pascal 下三角矩阵同型，其逆核为经典反二项变换：
\[
(L^{-1})_{a,b}=(-1)^{a+b-q}\binom{a}{q-b},
\]
代回即得所示闭式与整性。证毕。
\end{proof}

\begin{corollary}[例外谱倒数一次幂和的 Fibonacci 闭式]\label{cor:xi-terminal-replica-softcore-exceptional-reciprocal-sum-fibonacci}
对一切整数 \(q\ge 2\)，令 \(\{\lambda_i^{(q)}\}_{i=0}^{q}=\mathrm{spec}(A^{(q)})\)。
则成立严格恒等式
\[
\boxed{\ \sum_{i=0}^{q}\frac{1}{\lambda_i^{(q)}}
=\operatorname{tr}\bigl((A^{(q)})^{-1}\bigr)
=-1+2(-1)^qF_{q+1}\ } .
\]
\end{corollary}
\begin{proof}
由结论 \ref{con:xi-terminal-replica-softcore-Aq-explicit-inverse-integrality}，
\[
\operatorname{tr}\bigl((A^{(q)})^{-1}\bigr)
=\sum_{a=0}^{q}\bigl(2(-1)^q\binom{a}{q-a}-\delta_{a,0}\bigr)
=-1+2(-1)^q\sum_{a=0}^{q}\binom{a}{q-a}.
\]
令 \(k=q-a\) 并用 Fibonacci 的经典二项表示
\(
\sum_{k=0}^{\lfloor q/2\rfloor}\binom{q-k}{k}=F_{q+1}
\)
即得结论。证毕。
\end{proof}

\begin{corollary}[例外谱倒数二次幂和的 Fibonacci 闭式]\label{cor:xi-terminal-replica-softcore-exceptional-reciprocal-square-sum-fibonacci}
对一切整数 \(q\ge 2\)，令 \(\{\lambda_i^{(q)}\}_{i=0}^{q}=\mathrm{spec}(A^{(q)})\)。
则成立严格恒等式
\[
\boxed{\ \sum_{i=0}^{q}\frac{1}{\bigl(\lambda_i^{(q)}\bigr)^2}
:=\operatorname{tr}\bigl((A^{(q)})^{-2}\bigr)
=4F_{2q+2}+1\ } .
\]
\end{corollary}
\begin{proof}
由结论 \ref{con:xi-terminal-replica-softcore-Aq-explicit-inverse-integrality} 的证明，存在
\(
L_{a,b}=\binom{q-a}{b}
\)
与 \(e_0=(1,0,\dots,0)^{\mathsf T}\) 使
\[
(A^{(q)})^{-1}=2L^{-1}-e_0e_0^{\mathsf T}.
\]
因此
\[
\operatorname{tr}\bigl((A^{(q)})^{-2}\bigr)
=\operatorname{tr}\bigl((2L^{-1}-e_0e_0^{\mathsf T})^2\bigr)
=4\,\operatorname{tr}(L^{-2})-4\,e_0^{\mathsf T}L^{-1}e_0+\operatorname{tr}(e_0e_0^{\mathsf T}).
\]
由 \((L^{-1})_{0,0}=(-1)^{q}\binom{0}{q}=0\) 得 \(e_0^{\mathsf T}L^{-1}e_0=0\)，且 \(\operatorname{tr}(e_0e_0^{\mathsf T})=1\)，故
\[
\operatorname{tr}\bigl((A^{(q)})^{-2}\bigr)=4\,\operatorname{tr}(L^{-2})+1.
\]
另一方面，\(L=C^{(q)}\) 与 \(B_q^{\mathsf T}\) 相同（定理 \ref{thm:xi-replica-softcore-exceptional-interlacing} 的记号），而 \(B_q\) 的谱为
\(\{(-1)^k\varphi^{q-2k}\}_{k=0}^{q}\)（命题 \ref{prop:pom-bq-is-symq-and-spectrum}）。
故 \(L^{-2}\) 的谱为 \(\{\varphi^{-2(q-2k)}\}_{k=0}^{q}\)，从而
\[
\operatorname{tr}(L^{-2})
=\sum_{k=0}^{q}\varphi^{-2(q-2k)}
=\sum_{k=0}^{q}\varphi^{2(q-2k)}
=F_{2q+2},
\]
其中最后一步使用结论 \ref{con:xi-terminal-replica-softcore-exceptional-spectrum-square-sum} 证明中的等比求和恒等式。
代回即得所示闭式。证毕。
\end{proof}

\begin{corollary}[全矩阵行列式的塌缩与整数逆核]\label{cor:xi-terminal-replica-softcore-Tq-det-and-integer-inverse-kernel}
设 \(q\ge 2\) 并记 \(N:=2^q\)。
则
\[
\boxed{\ \det T^{(q)}=2^{-(N-1)}\ }.
\]
并且 \(\bigl(T^{(q)}\bigr)^{-1}\) 为整数矩阵：对任意 \(u,v\in U_q\)，记 \(\mathbf{0}:=(0,\dots,0)\)、\(\mathbf{1}:=(1,\dots,1)\)，则
\[
\boxed{\ \bigl(T^{(q)}\bigr)^{-1}_{u,v}
=2(-1)^{|u\wedge v|}\,\mathbf{1}_{\{u\vee v=\mathbf{1}\}}
-\delta_{u,\mathbf{0}}\delta_{v,\mathbf{0}}\ } .
\]
\end{corollary}
\begin{proof}
由 \(B^{(q)}:=2T^{(q)}\) 与结论 \ref{con:xi-terminal-disjointness-bq-smith-snf-coker-parity} 的 \(\det(B^{(q)})=2\)，
\(
\det(T^{(q)})=2^{-N}\det(B^{(q)})=2^{-(N-1)}
\)。
另一方面，结论 \ref{con:xi-terminal-disjointness-bq-complement-selfdual-inverse-rankone} 给出
\(
(B^{(q)})^{-1}=(D^{(q)})^{-1}-\frac12E_{\varnothing,\varnothing}
\)
且 \((D^{(q)})^{-1}_{u,v}=(-1)^{|u\cap v|}\mathbf{1}_{\{u\cup v=[q]\}}\)；
将 \(U_q\cong\mathcal P([q])\) 复原为按位 \(\wedge,\vee\) 记号并用 \(T^{(q)}=\frac12B^{(q)}\) 即得所示整数核。证毕。
\end{proof}

\begin{corollary}[两层核的行列式闭式：仅由 \texorpdfstring{$(a,b,N)$}{(a,b,N)} 决定]\label{cor:xi-terminal-replica-softcore-twolevel-kernel-det}
设 \(q\ge 2\) 并记 \(N:=2^q\)。
对任意标量 \(a,b\in\CC\)，定义两层核矩阵 \(T^{(q)}(a,b)\in\CC^{U_q\times U_q}\) 为
\[
T^{(q)}(a,b)_{u,v}:=
\begin{cases}
a,&u\wedge v=\mathbf{0},\\[2pt]
b,&u\wedge v\neq \mathbf{0}.
\end{cases}
\]
则成立严格闭式
\[
\boxed{\ \det\bigl(T^{(q)}(a,b)\bigr)=a\,(a-b)^{N-1}\ }.
\]
\end{corollary}
\begin{proof}
当 \(a=b\) 时，\(T^{(q)}(a,a)=a\,\mathbf{1}\mathbf{1}^{\mathsf T}\) 为秩一矩阵且 \(N\ge 4\)，故行列式为 \(0\)，与右端一致。
以下设 \(a\neq b\)。
\par
由定义
\(
T^{(q)}(a,b)=b\,\mathbf{1}\mathbf{1}^{\mathsf T}+(a-b)D^{(q)}
\)。
由于 \(D^{(q)}=K^{\otimes q}\) 且 \(q\ge 2\)，有 \(\det(D^{(q)})=1\)（见结论 \ref{con:xi-terminal-disjointness-bq-smith-snf-coker-parity} 的证明）。
并且由 \(\bigl(D^{(q)}\bigr)^{-1}\mathbf{1}=e_{\varnothing}\) 得
\(\mathbf{1}^{\mathsf T}\bigl(D^{(q)}\bigr)^{-1}\mathbf{1}=1\)。
对矩阵行列式引理应用于
\(
A:=(a-b)D^{(q)}
\)
与秩一更新 \(b\,\mathbf{1}\mathbf{1}^{\mathsf T}\)，得到
\[
\det\bigl(T^{(q)}(a,b)\bigr)
=\det(A)\Bigl(1+b\,\mathbf{1}^{\mathsf T}A^{-1}\mathbf{1}\Bigr)
=(a-b)^{N}\det(D^{(q)})\Bigl(1+\frac{b}{a-b}\Bigr)
=a\,(a-b)^{N-1}.
\]
证毕。
\end{proof}

\begin{corollary}[逆核支撑的精确计数：\texorpdfstring{$3^q+1$}{3^q+1} 稀疏度与逐行分布]\label{cor:xi-terminal-replica-softcore-inverse-support-count}
沿用推论 \ref{cor:xi-terminal-replica-softcore-Tq-det-and-integer-inverse-kernel} 的显式逆核。
则 \(\bigl(T^{(q)}\bigr)^{-1}\) 的非零元总数满足
\[
\#\Bigl\{(u,v)\in U_q^2:\ \bigl(T^{(q)}\bigr)^{-1}_{u,v}\neq 0\Bigr\}=3^q+1.
\]
更精确地，若 \(u\neq \mathbf{0}\)，则第 \(u\) 行恰有 \(2^{|u|}\) 个非零元，且其取值集合包含于 \(\{\pm 2\}\)；
而第 \(\mathbf{0}\) 行恰有两个非零元：
\[
\bigl(T^{(q)}\bigr)^{-1}_{\mathbf{0},\mathbf{0}}=-1,
\qquad
\bigl(T^{(q)}\bigr)^{-1}_{\mathbf{0},\mathbf{1}}=2.
\]
\end{corollary}
\begin{proof}
由推论 \ref{cor:xi-terminal-replica-softcore-Tq-det-and-integer-inverse-kernel}，
当 \((u,v)\neq(\mathbf{0},\mathbf{0})\) 时有
\(
\bigl(T^{(q)}\bigr)^{-1}_{u,v}\neq 0
\iff
u\vee v=\mathbf{1}
\)，且此时取值必为 \(\pm 2\)。
因此对固定 \(u\neq \mathbf{0}\)，条件 \(u\vee v=\mathbf{1}\) 要求 \(v\) 在 \(\{j:\ u_j=0\}\) 上强制取 \(1\)，并在 \(\{j:\ u_j=1\}\) 上自由取 \(0/1\)，从而解的个数为 \(2^{|u|}\)。
对 \(u=\mathbf{0}\) 时仅有 \(v=\mathbf{1}\) 满足 \(u\vee v=\mathbf{1}\)，且 \((\mathbf{0},\mathbf{0})\) 处另有修正项 \(-1\)，得到所示两处非零值。
\par
将 \(2^{|u|}\) 对所有 \(u\in U_q\) 求和得
\[
\sum_{u\in U_q}2^{|u|}
=\sum_{k=0}^{q}\binom{q}{k}2^{k}
=(1+2)^q
=3^q,
\]
对应于 \(u\vee v=\mathbf{1}\) 产生的全部非零元；再加上额外的 \((\mathbf{0},\mathbf{0})\) 条目即得总数 \(3^q+1\)。
证毕。
\end{proof}

\begin{conclusion}[Fibonacci 幂矩恒等式与 Perron 根的收敛级数方程]\label{con:xi-terminal-replica-softcore-fib-moment-and-perron-series}
沿用结论 \ref{con:xi-terminal-replica-softcore-fiber-splitting-full-spectrum} 的记号。
对任意整数 \(m\ge 0\)，成立严格幂矩恒等式
\[
\boxed{\ \mathbf{1}^{\mathsf T}(D^{(q)})^{m}\mathbf{1}=F_{m+3}^{\,q}\ }.
\]
记 Perron 根（最大特征值）为 \(\rho(T^{(q)})\)。
则当 \(\lambda>\varphi^{q}/2\) 时有收敛的 resolvent 展开，并且 \(\lambda=\rho(T^{(q)})\) 满足严格方程
\[
\boxed{\ 1=\frac12\sum_{m\ge 0}\frac{F_{m+3}^{\,q}}{2^{m}\lambda^{m+1}}\ } .
\]
\end{conclusion}
\begin{proof}
由 \(D^{(q)}=K^{\otimes q}\) 与 \(\mathbf{1}=u^{\otimes q}\)（\(u=(1,1)^{\mathsf T}\)），
\[
\mathbf{1}^{\mathsf T}(D^{(q)})^{m}\mathbf{1}
=\bigl(u^{\mathsf T}K^{m}u\bigr)^q.
\]
而 \(K\) 为 Fibonacci \(Q\)-矩阵，满足
\(
K^{m}=\bigl(\begin{smallmatrix}F_{m+1}&F_m\\ F_m&F_{m-1}\end{smallmatrix}\bigr)
\)；
故 \(u^{\mathsf T}K^{m}u=F_{m+1}+2F_m+F_{m-1}=F_{m+3}\)，得到幂矩恒等式。
\par
对 Perron 根，取正特征向量 \(x\) 并令 \(\mathbf{1}^{\mathsf T}x\neq 0\)。
由
\(
(\lambda I-\tfrac12D^{(q)})x=\tfrac12(\mathbf{1}^{\mathsf T}x)\mathbf{1}
\)
消去 \(\mathbf{1}^{\mathsf T}x\) 得
\(
1=\tfrac12\,\mathbf{1}^{\mathsf T}(\lambda I-\tfrac12D^{(q)})^{-1}\mathbf{1}
\)。
当 \(\lambda>\|D^{(q)}/2\|=\varphi^{q}/2\) 时，Neumann 级数
\[
(\lambda I-\tfrac12D^{(q)})^{-1}
=\sum_{m\ge 0}\frac{(D^{(q)})^m}{2^{m}\lambda^{m+1}}
\]
收敛；代入并用幂矩恒等式即得所示级数方程。证毕。
\end{proof}

\endinput

